\section*{Chapter \ref{ch:b3:competition-policy} --- Competition Policy and Regulation}

\begin{solution}{exo:b3:competition-policy:1}
By $3.4 - 2.2 = 1.2$ points a year. Traffic growth slowed, from 9.1 to 8.2~\% a year, and to 3.5~\% in the
following two decades. % figdata: yields
\end{solution}

\begin{solution}{exo:b3:competition-policy:2}
$900 + 625 + 400 + 225 + 100 = 2{,}250$. The merger adds $2 \times 15 \times 10 = 300$: 2{,}550. Above 1{,}800 with a change above 100: presumed
to lessen competition substantially.
\end{solution}

\begin{solution}{exo:b3:competition-policy:3}
Separate: $w = 60$, $P = 80$, $Q = 20$. Integrated: $P = 60$, $Q = 40$.
\end{solution}

\begin{solution}{exo:b3:competition-policy:4}
Retail prices that rise after the restraint is adopted while the quantities sold
fall, with no gain in measured service; under the service story quantities rise.
\end{solution}

\begin{solution}{exo:b3:competition-policy:5}
$0.05/(0.05 + 0.40) = 11.1~\%$.
\end{solution}

\begin{solution}{exo:b3:competition-policy:6}
$0.225/2 = 0.1125$: prices are on average 11.25~\% above marginal cost.
\end{solution}

\begin{solution}{exo:b3:competition-policy:7}
$10 - 0.5 \times 2/0.5 = 8$. With $Q = \sqrt{KL}$ cost minimisation sets $K/L$ equal to the ratio of input prices, $w$ over
the capital price: $10/8 = 1.25$, against $10/10 = 1$.
\end{solution}

\begin{solution}{exo:b3:competition-policy:8}
The tested solution is \code{ex08_hhi.py}, in the chapter's \code{exercises/}
folder. Cement: 2{,}250 to 2{,}550, presumed. Flour: 2{,}000 to 2{,}800, presumed. Paint: 2{,}944 to 3{,}184, presumed although
the merged share is only 22~\%. Glass: 1{,}900 to 1{,}950, a change of 50, not presumed although the
market is highly concentrated. % output: ex08_hhi
\end{solution}

\begin{solution}{exo:b3:competition-policy:9}
Fares on the route falling no faster than the trend of the regulated years, or no
lower than a cost-based counterfactual such as the regulator's formula, once
changes in costs and quality are allowed for --- or fares rising where one airline
came to dominate the route's airports.
\end{solution}

\begin{solution}{exo:b3:competition-policy:10}
The proof is in the text. The formula fixes the average markup over each firm's own
cost. If the merged firm's cost falls, its output and share rise and the HHI with them,
but the price can still fall, because the costs the markups are added to are lower.
\end{solution}

\begin{solution}{exo:b3:competition-policy:11}
The tested solution is \code{ex11_cap.py}, in the chapter's \code{exercises/}
folder. The cap falls from 99.00 to 95.10 while unit cost falls from 86.40 to 73.38; profit rises from 12{,}600
to 21{,}715. Under rate of return the year-5 price is 83.38 and profit 10{,}000. % output: ex11_cap
The cap gives consumers less at first and the firm a reason to cut costs; at the next
review the regulator can pass the lower costs on.
\end{solution}

\begin{solution}{exo:b3:competition-policy:12}
Condemn when the probability of harm exceeds $1/(1 + 3) = 0.25$. The prior odds are $0.2/0.8 = 0.25$; doubled, 0.5, a
probability of $0.5/1.5 = 0.33$: condemn.
\end{solution}

\begin{solution}{pb:b3:competition-policy:water}
% figdata: water_stats
\textbf{1.} Minimising $10L + 10K$ subject to $\sqrt{KL} = Q$ gives $K = L = Q$ and cost $20Q$: marginal cost 20.

\textbf{2.} Average cost $20 + 600/Q$ falls with output over the whole range of demand: one network
is cheapest.

\textbf{3.} $P = 60$, $Q = 40$, profit $40 \times 40 - 600 = 1{,}000$.

\textbf{4.} $40^2/2 = 800$.

\textbf{5.} $\tfrac12 \times (80 - 40) \times (60 - 20) = 800$.

\textbf{6.} $Q = 80$, consumer surplus 3{,}200, a loss of 600.

\textbf{7.} $(100 - Q)Q = 20Q + 600$ gives $Q = 40 + \sqrt{1{,}000} = 71.6$ and $P = 28.38$.

\textbf{8.} $71.62^2/2 = 2{,}564.9$; deadweight loss $\tfrac12 \times (80 - 71.62) \times (28.38 - 20) = 35.1$.

\textbf{9.} At marginal cost, taxpayers, through a subsidy of 600; at average cost, water users.

\textbf{10.} $P = 40$; $K = L = 60$, cost 1{,}200.

\textbf{11.} $40 \times 60 - 1{,}200 - 600 = 600$; allowed $(12 - 10) \times 60 = 120$.

\textbf{12.} With the constraint binding, revenue less labour and fixed costs is $12K$, so profit is
$12K - 10K = 2K$.

\textbf{13.} The firm takes the largest $K$ with $2{,}400 - 600 = 10 \times 3{,}600/K + 12K$: $K = 126.2$, $L = 28.5$, a ratio of 4.43 against 1.

\textbf{14.} 1{,}547.5 with the unrounded inputs ($10 \times 28.52 + 10 \times 126.23$), an excess of 347.5 over 1{,}200 (29~\%).

\textbf{15.} It observes costs and capital, not the cost-minimising plan; it would need the
production function to know the minimum.

\textbf{16.} $(28.38 - 16) \times 71.62 - 600 - 150 = 136.5$, against 0 without the effort: it makes it.

\textbf{17.} $(100 - Q)Q = 16Q + 600$ gives $Q = 76.1$ and $P = 23.88$; profit $0 - 150 = -150$: it does not.

\textbf{18.} $X = 1.03$ dalens.

\textbf{19.} $P = 27.35$, consumer surplus 2{,}639.3 (a gain of 74.4) and profit 74.4.

\textbf{20.} The effort is paid before its gains; a firm that expects $X$ to be raised
once costs fall gains less from the effort and may skip it.

\textbf{21.} The price falls to 23.88 and the firm's gain disappears; firms that foresee
this ratchet cut costs less.

\textbf{22.} If its price follows the costs of comparable firms, its own savings do not
lower its own price, so it keeps them, while consumers gain as the yardstick falls.

\textbf{23.} Monopoly 800, average cost 2{,}564.9, the cap with $X = 1.03$ 2{,}639.3.

\textbf{24.} Marginal-cost pricing, 3{,}200; it needs a subsidy of 600 a day from taxes, which have
their own cost (\cref{ch:b3:efficiency-cost-of-taxation}), and it gives the firm no reason to cut costs.

\textbf{25.} Consumer surplus and the firm's profit are 800 and 1{,}000 under monopoly, 3{,}200 and $-600$ at
marginal cost, 2{,}564.9 and 0 at average cost, and 2{,}639.3 and 74.4 under the cap with $X = 1.03$, the offset
that shares the cost saving equally.
\end{solution}
